% preambulo.tex — cada \usepackage comentado com a regra do guia que atende.
% Fonte de todas as remissões: auditoria/parametros_guia.md (página impressa do guia).
%
% Decisão de classe (Fase 2, registrada em auditoria/log.md): abntex2, não formatação
% manual. Justificativa: abntex2 já nasce configurado para as margens, corpo de fonte
% e entrelinhas que o guia pede (§1 de parametros_guia.md — 3/2cm, 12pt, 1,5 entrelinhas),
% e resolve de fábrica a maior parte dos elementos pré-textuais (capa, folha de rosto,
% ficha catalográfica, resumo, listas, sumário — §7). Onde o guia diverge do
% comportamento padrão do pacote, o pacote é sobrescrito e a divergência vai para o log
% (R-3: guia vence pacote). Montar tudo à mão custaria mais do que essas sobrescritas.

\documentclass[
  12pt,          % corpo do texto — guia p. 12 ("fonte tamanho 12 para o texto")
  a4paper,       % papel A4 (210x297mm) — guia p. 12
  oneside,       % entrega só digital (PDF); guia permite anverso-only para textual/
                 % pós-textual (p. 13) e não exige twoside para TCC de graduação.
                 % Único caso de "verso" no guia é a ficha catalográfica atrás da folha
                 % de rosto (p. 13), tratado como página própria em sequência, não como
                 % impressão física frente/verso — ver DECISÃO em auditoria/log.md.
  chapter=TITLE, % títulos de capítulo em caixa alta + negrito — confirmado no exemplo
                 % real de sumário do próprio guia (Figura 11, p. 20: "2 A BIBLIOTECA"
                 % em maiúsculas e negrito). NÃO usar section=TITLE: no mesmo exemplo,
                 % os títulos de nível 2 ("2.1 História") aparecem em negrito, mas
                 % SEM caixa alta. O negrito em si NÃO é automático (achado C2/Fase 8,
                 % ver bloco de correção logo abaixo do \documentclass) — a suposição
                 % original aqui estava errada.
  english,       % hifenização auxiliar (títulos/termos em inglês no corpo)
  brazil,        % idioma principal do documento
]{abntex2}

% memhfixc: abntex2 é baseado em memoir, que já carrega hyperref internamente;
% memoir+hyperref têm uma interação conhecida de contadores/âncoras que produz
% destinos de página duplicados (ex. "page.1" repetido entre capa e folha de
% rosto) — memhfixc é o patch oficial do próprio memoir para isso, e precisa vir
% DEPOIS de hyperref já estar carregado (por isso aqui, logo após \documentclass,
% já que abntex2.cls carrega hyperref internamente antes deste ponto).
\usepackage{memhfixc}

% Negrito em título de capítulo e de seção primária (guia p. 21-22 e Figura 11,
% p. impressa 20) — achado C2/Fase 8 (ver auditoria/log.md §21): abntex2.cls
% define \ABNTEXchapterfont e \ABNTEXsectionfont como \sffamily puro, sem
% \bfseries em nenhum dos dois, e \setsecheadstyle/\chaptitlefont usam essas
% macros diretamente — nenhum nível de título sai em negrito por padrão. Só a
% capa (\imprimirtitulo) tem \bfseries embutido manualmente na classe, prova de
% que o negrito é intencional em outros títulos de destaque, só ausente aqui.
% \ABNTEXsubsectionfont/\ABNTEXsubsubsectionfont, na classe, são definidas como
% referências a \ABNTEXsectionfont (não cópias) — por isso são redefinidas aqui
% também, de forma independente, sem \bfseries, para NÃO herdar o negrito que
% \ABNTEXsectionfont ganha logo abaixo (guia: só seção 1ª/2ª em negrito).
\renewcommand{\ABNTEXchapterfont}{\sffamily\bfseries}
\renewcommand{\ABNTEXsectionfont}{\sffamily\bfseries}
\renewcommand{\ABNTEXsubsectionfont}{\sffamily}
\renewcommand{\ABNTEXsubsubsectionfont}{\sffamily}

% Mesmo defeito, mesma causa, achado no recheque do fix acima (Fase 8, ver
% auditoria/log.md §21): no Sumário, \cftchapterfont/\cftsectionfont
% (abntex2.cls linhas 520-529) são \bfseries, mas \cftsubsectionfont/
% \cftsubsubsectionfont só ajustam \normalsize/\small, sem nunca desligar o
% negrito com \mdseries — confirmado por `pdffonts`: a página do Sumário só
% usa a fonte Bold, em todo nível, do capítulo à subsubseção. Corrigido do
% mesmo jeito: \mdseries explícito só nesses dois níveis, sem tocar
% \cftchapterfont/\cftsectionfont (que devem continuar em negrito).
\renewcommand{\cftsubsectionfont}{\normalsize\mdseries}
\renewcommand{\cftsubsubsectionfont}{\small\mdseries}

% ---------------------------------------------------------------------------
% Idioma e codificação
% ---------------------------------------------------------------------------
\usepackage[T1]{fontenc}
\usepackage[utf8]{inputenc}
\usepackage{lmodern}         % fallback de fonte vetorial; substituído por mathptmx abaixo
\usepackage{mathptmx}        % Times New Roman (equivalente PostScript) — guia p. 12
                              % permite Arial OU Times New Roman; Times é o padrão de
                              % fábrica do ecossistema abntex2/ABNT em pdfLaTeX.

% ---------------------------------------------------------------------------
% Citação (NBR 10520:2023) e referências (NBR 6023:2025)
% ---------------------------------------------------------------------------
\usepackage[
  alf,
  abnt-etal-text=it,         % *et al.* em itálico — regra já vigente no texto-fonte
  abnt-doi=doi,              % NÃO expande "Disponível em:" para um link
                              % https://dx.doi.org/... — testado com
                              % abnt-doi=expand e não teve efeito nenhum: o
                              % ramo de compose.url() que consultaria essa
                              % opção só roda com abnt.url.package=0, e este
                              % documento usa abnt.url.package=1 (pacote url
                              % ativo, citeoption "abnt-url-package=url" já
                              % disparada pelo próprio abntex2cite), cujo
                              % ramo ignora abnt-doi por completo — achado
                              % empírico, não vale a pena outro patch de .bst
                              % só por isso. A cláusula "DOI: 10.xxxx."
                              % separada que o guia pede (p. 34: "...ano. DOI
                              % (se houver). Disponível em:...") NÃO existe
                              % nesta versão do abntex2-alf.bst nem como
                              % opção — patch local em
                              % latex/bst/abntex2-alf-tg1.bst, função
                              % format.url (achado do autor, ver
                              % auditoria/log.md §25). "Disponível em:"
                              % continua mostrando "doi:10.xxxx" cru, o que
                              % não é vedado pelo guia (nota técnica de
                              % parametros_guia.md §9.1: os sinais < > são
                              % convenção tipográfica, não texto do URL).
  abnt-full-initials=yes,   % prenomes por extenso na lista — regra do projeto
                              % ("Prenomes por extenso em todas as entradas").
                              % Default do pacote é abreviar (achado empírico:
                              % 1ª compilação saiu "ACHTERBERG, T." em vez de
                              % "ACHTERBERG, Tobias" — ver auditoria/log.md).
  tg1-cite-style,            % "(Author, YEAR)" — inicial maiúscula na chamada,
                              % caixa alta só na lista (NBR 10520:2023, já em
                              % vigor nos dois capítulos-fonte). Chave própria,
                              % definida em referencias.bib, citando um valor
                              % que o pacote embutido também define, mas sob uma
                              % chave que colide (achado empírico, auditoria/
                              % log.md Fase 4: BibTeX trata chaves de citação
                              % como insensíveis a maiúscula/minúscula, e
                              % "abnt-cite-style=AUTHORYEAR" colide com
                              % "...=AuthorYEAR", sempre resolvendo para a
                              % primeira, caixa alta — daí a chave própria).
                              % Passada como OPÇÃO DE PACOTE, não via
                              % \citeoption direto: \citeoption chamado no
                              % preâmbulo escreve no .aux antes de \@auxout
                              % existir e é descartado em silêncio; como opção
                              % de pacote, o próprio abntex2cite adia a chamada
                              % para \AtBeginDocument, quando o .aux já existe.
]{abntex2cite}
% Estilo alfabético (autor-data) — catálogo já usa esse sistema em todo o texto-fonte
% (ex. "(FORREST et al., 2019)" na lista, "(Forrest et al., 2019)" em texto), não o
% sistema numérico.

% ---------------------------------------------------------------------------
% Figuras, quadros e tabelas — legenda acima, fonte abaixo (guia p. 23-24)
% ---------------------------------------------------------------------------
\usepackage{graphicx}
\usepackage{float}
\usepackage{caption}
\usepackage{placeins}       % \FloatBarrier — achado do autor, 13/08/2026 (ver
                              % auditoria/log.md §25): um Quadro[htbp] adiado
                              % (Quadro 7) drenava a página seguinte inteira e
                              % empurrava a Figura 3 (mesmo com [H]) 2 páginas
                              % para frente. \FloatBarrier força os flutuantes
                              % pendentes de uma seção a resolver antes da
                              % próxima, sem alterar nenhuma legenda/numeração.

% Legenda ACIMA do conteúdo para figuras — diverge do padrão de fábrica do LaTeX
% (que põe legenda de figura abaixo). DIVERGÊNCIA guia×pacote, registrada em
% auditoria/log.md: o guia (p. 23) manda identificação "na parte superior" para
% qualquer tipo de ilustração, sem exceção para figuras.
\floatstyle{plaintop}
\restylefloat{figure}
\restylefloat{table}

% Ambiente \quadro distinto de \table — plano exige que a diferença entre Quadro
% (qualitativo) e Tabela (numérico) seja estrutural, não editorial (prompt, §5).
% Numeração corrida (sem "within"), não por capítulo: o texto-fonte já audita
% Quadro 1..11 (mais 3a e 7a) como sequência única do documento, não "Quadro
% 3.1" — ver auditoria/log.md, Fase 3.
\newfloat{quadro}{htbp}{loq}
\floatname{quadro}{Quadro}
\floatstyle{plaintop}
\restylefloat{quadro}

% Figura e Tabela também numeradas de forma corrida, pelo mesmo motivo: o
% texto-fonte usa "Figura 1/2/3" e "Tabela 1" flat, não escopado por capítulo
% (que é o padrão de fábrica do memoir/abntex2 para relatórios multi-capítulo).
\counterwithout{figure}{chapter}
\counterwithout{table}{chapter}
% Equações também: numeração corrida, não "(3.1)" por capítulo. Só as 2
% equações efetivamente referenciadas em texto (PL, MILP) são numeradas — as
% outras 5 do Cap. 3 usam align*/equation* (guia p. 22: numera-se "só quando
% necessário"; achado C2/Fase 8, ver auditoria/log.md §19).
\counterwithout{equation}{chapter}

% "Título de seção primária inicia em página ímpar" (guia p. 12) é regra de
% NUMERAÇÃO, independente de oneside/twoside (que é só sobre impressão física).
% A opção nativa \openright da memoir NÃO resolve isso aqui: seu
% \clearforchapter (memoir.cls) só insere a página em branco sob \if@twoside —
% em oneside ele degrada para um \clearpage comum, achado empírico (Fase 8).
% Reimplementado manualmente, sem depender de \if@twoside: mantém oneside
% (decisão da Fase 2, entrega só digital) e ainda força \chapter a cair em
% página ímpar. Achado C2/Fase 8, ver auditoria/log.md §19.
\makeatletter
\renewcommand{\clearforchapter}{%
  \clearpage
  \ifodd\value{page}\else
    \hbox{}\thispagestyle{empty}\newpage
  \fi
}
\makeatother

% Comando \fontenota — a "Fonte:" abaixo do conteúdo é elemento obrigatório mesmo em
% produção do próprio autor (guia p. 23: "elaborado pelo próprio autor"). Não é o
% \caption padrão do LaTeX (que abntex2 já usa para a legenda de cima); é uma segunda
% linha, centralizada, fonte reduzida, adicionada manualmente dentro de cada float.
\newcommand{\fontenota}[1]{%
  \par\nobreak
  {\small\centering Fonte: #1\par}%
}
% \par antes do grupo: sem isso, uma tabela/quadro estreito (que não ocupa a
% largura inteira, ex. a Tabela 1 de McNemar) não força quebra de linha
% sozinho, e a linha "Fonte:" aparece colada ao lado do conteúdo em vez de
% abaixo — achado empírico, corrigido aqui de uma vez para todos os usos.

\captionsetup{
  labelsep=endash,   % "travessão" entre número e título — guia p. 23 ("travessão")
  font=small,        % fonte reduzida para legendas — guia p. 12 (10/11 para legendas)
  justification=centering,
}

% ---------------------------------------------------------------------------
% Equações numeradas à direita, só as referenciadas — guia p. 22
% ---------------------------------------------------------------------------
\usepackage{amsmath}
\usepackage{amssymb}
% amsmath numera à direita por padrão; equações sem \label/\eqref não citadas
% permanecem sem número visível de fato relevante — a disciplina de "só numerar o
% que é referenciado" é aplicada na redação (Fase 3), não no preâmbulo.

% ---------------------------------------------------------------------------
% Apoio geral
% ---------------------------------------------------------------------------
\usepackage{indentfirst}
% Recuo de primeira linha de parágrafo — o guia NÃO especifica um valor de recuo em
% nenhuma das 66 páginas lidas (só especifica o recuo de 4cm para citação longa,
% guia p. 25). DECISÃO (não DIVERGÊNCIA, porque o guia é silente, não contraditório):
% usar o recuo padrão do abntex2 (1,25cm), prática corrente ABNT. Ver auditoria/log.md.

\usepackage{microtype}       % justificação tipográfica fina — não normativo, só qualidade
\usepackage{xcolor}          % suporte de cor para BPMN e destaques pontuais
\usepackage{booktabs}        % regras horizontais para a Tabela de McNemar (única Tabela)
\usepackage{tabularx}        % colunas de largura flexível — os 13 Quadros do Cap. 3
                              % têm células de texto longo, não cabem em tabular fixo
\usepackage{array}           % suporte a especificadores de coluna customizados (p{})

% hyperref já é carregado internamente pelo abntex2 (via bookmark.sty) — carregá-lo
% de novo com opções próprias causa "Option clash". Ajuste feito por \hypersetup,
% não por novo \usepackage. Não é exigido pelo guia; só navegação interna do PDF
% (sumário clicável, remissões), sem alterar a aparência impressa do documento.
% plainpages=false: sem isso, hyperref nomeia âncoras de página pelo valor impresso
% de \thepage, que o abntex2 mantém intencionalmente parado durante \pretextual
% (guia p. 12: "folhas pré-textuais devem ser contadas... mas não numeradas") —
% duas páginas físicas distintas (capa, folha de rosto) ficam com o mesmo \thepage
% e colidem no nome da âncora ("destination... page.1... duplicate ignored").
% plainpages=false usa o número físico absoluto para o nome da âncora em vez do
% valor impresso, resolvendo a colisão sem alterar a numeração visível do documento.
% Fonte do diagnóstico: TeX FAQ, "pdfTeX destination ... ignored"
% (https://texfaq.org/FAQ-hyperdupdest).
\hypersetup{colorlinks=false, hidelinks, plainpages=false, pageanchor=false}

% ---------------------------------------------------------------------------
% Metadados do documento — preenchidos em pretextual/dados.tex (Fase 7)
% ---------------------------------------------------------------------------
% \autor, \titulo, \instituicao etc. ficam num arquivo à parte porque dependem de
% dados que ainda não foram confirmados (nome completo do autor, título exato,
% orientador) — ver PENDÊNCIA em auditoria/log.md. Não fabricar aqui (R-2).
